\section{The Migdal Effect}
\label{sec:migdal}

\subsection{Physical process: why it is needed}

For elastic DM--nucleus scattering, the maximum nuclear recoil energy is
$E_n^\mathrm{max} = 2\munuc^2 v^2/m_N$. For a Si nucleus ($m_N\simeq26$~GeV)
and $v\sim10^{-3}c$:

\begin{center}
\begin{tabular}{@{}lll@{}}
\toprule
$m_\chidm$ & $E_n^\mathrm{max}$ & Detectable in a CCD? \\
\midrule
1~GeV & $\sim$100~eV & Yes \\
100~MeV & $\sim$1~eV & Marginal \\
10~MeV & $\sim$0.01~eV & No \\
\bottomrule
\end{tabular}
\end{center}

Below $\sim$100~MeV, the nuclear recoil alone is invisible. The
\textbf{Migdal effect} converts part of the recoil momentum into an
\emph{electronic} excitation via a sudden-perturbation (shake-off) process,
which \emph{can} be detected as an ionization signal — extending
sensitivity down to $\sim$1~MeV.

\subsection{The sudden approximation and factorization}

Because the nuclear scattering timescale ($\sim1/q\sim10^{-22}$~s) is far
shorter than the electronic relaxation timescale
($\sim1/\omega\sim10^{-17}$--$10^{-15}$~s), the process factorizes into two
sequential steps:
\[
\underbrace{\chidm + \text{nucleus}_A \to \chidm' + \text{nucleus}_A^*}_{\text{nuclear scattering}}
\qquad\longrightarrow\qquad
\underbrace{\text{nucleus}_A^* \to \text{nucleus}_A + e^- + \ldots}_{\text{electronic shake-off}}
\]
The differential rate for an electronic excitation of energy $\omega$ is
\begin{equation}
\frac{dR}{d\omega} = \frac{\rhochi}{m_N m_\chidm}\, I(\omega)
\int_{\vmin}^{\vesc+v_E} v\,f(v)\,J(v,\omega)\, dv\,.
\end{equation}

\paragraph{$I(\omega)$ — shake-off probability (material physics):}
\begin{equation}
I(\omega) = \frac{1}{E_n}\frac{dP}{d\omega}
= \frac{2\alpha_\mathrm{EM}\,Z_\mathrm{ion}^2}{3\pi^2\,\omega^4}
\int_0^{k_\mathrm{max}} k^2\,\Im\!\left[-\frac{1}{\varepsilon(\omega,k)}\right] dk\,,
\end{equation}
driven by the \textbf{same ELF}, $\Im[-1/\varepsilon(\omega,k)]$, that
governs DM--electron scattering (Section~\ref{sec:dme}) and dark photon
absorption (Section~\ref{sec:darkphoton}) — the Migdal effect is literally
the electron energy-loss function of the material, probed by a recoiling
nucleus instead of an external probe.

\paragraph{$J(v,\omega)$ — nuclear kinematics (DM/mediator physics),
free-nucleus approximation:}
\begin{equation}
J(v,\omega) \propto \frac{A^2\,\sigman}{v\,\munuc^2}\,\FDM^2(q)\,,
\end{equation}
with $A^2=784$ for Si ($A=28$) the coherent nuclear enhancement, and
$\sigman$ the DM-nucleon reference cross section at $q_0=\alpha m_e$.

\paragraph{Kinematic threshold:}
\begin{equation}
\vmin(\omega) = \sqrt{2\omega/\munuc}\,.
\end{equation}
For $\omega\simeq\Egap=1.11$~eV and $m_\chidm=10$~MeV, $\vmin\approx6\times10^{-3}c$
— already strongly suppressed by the halo tail, setting the effective
low-mass floor of the Migdal search.

\subsection{Mediator form factor (DM--nucleus vertex)}

\begin{itemize}[leftmargin=*]
  \item \textbf{Heavy mediator}: $\FDM^2(q) = (q_0^2+m_\mathrm{med}^2)^2/(q^2+m_\mathrm{med}^2)^2
  \to 1$ (thermal freeze-out benchmark, $q$-independent).
  \item \textbf{Light (massless) mediator}: $\FDM^2(q) = (q_0/q)^4$
  (freeze-in benchmark, strongly enhanced at low $q$).
\end{itemize}
This is a \textbf{distinct namespace} from the DM--electron mediator of
Section~\ref{sec:dme} — same ``heavy/light'' language, different vertex
(DM--nucleus vs.\ DM--electron).

\subsection{Standard silicon case}

\begin{center}
\begin{tabular}{@{}p{4.3cm}p{4cm}p{5.5cm}@{}}
\toprule
Parameter & Value & Source \\
\midrule
$\Egap$ & 1.11~eV & darkelf Si optical data, 130~K \\
$\bar\omega$ (phonon frequency) & 0.03~eV & Si \\
$E_\mathrm{nth}$ (nuclear recoil cutoff) & $0.12$~eV $=4\bar\omega$ &
  free-nucleus approximation validity floor \\
$Z_\mathrm{ion}(0)$ & 4.0 & effective ionic charge (\texttt{Si\_Zion.dat},
  \citealt{Brown2006}) \\
$A$ & 28 & Si mass number $\to A^2=784$ coherent enhancement \\
ELF source & \texttt{Si\_mermin.dat} (Mermin extension of optical data) &
  darkelf \\
Approximation & free-nucleus, \texttt{method="grid"} & impulse
  approximation and Ibe atomic Migdal deferred \\
SHM & $v_0=238$, $v_E=253.7$, $\vesc=544$~km/s & darkelf default \\
\bottomrule
\end{tabular}
\end{center}

\paragraph{Reference rate} (heavy mediator, $\sigman=10^{-36}$~cm$^2$,
integrated $\omega\in[\Egap,20\,\mathrm{eV}]$):

\begin{center}
\begin{tabular}{@{}lc@{}}
\toprule
$m_\chidm$ & Rate [events/kg/yr] \\
\midrule
100~MeV & $\approx3.5\times10^4$ \\
1~GeV & $\approx1.1\times10^6$ \\
\bottomrule
\end{tabular}
\end{center}

Rate scales \textbf{exactly linearly} in $\sigman$ (it enters only as an
overall prefactor of $J$), which the grid generator exploits as a fast path.
The DAMIC-M PRL 2025 reports the Migdal channel (using darkelf) as the
\textbf{most stringent limit for $m_\chidm=1$--$35$~MeV/$c^2$}, with the
explicit caveat that ``the Migdal effect is not yet experimentally
calibrated'' \citep{DAMICM2025}.

The detector pipeline from \texttt{ChargeIonization} onward is
\textbf{100\% identical} to DM--electron scattering: darkelf's
$dR_\mathrm{Migdal}/d\omega(\omega)$ already outputs the spectrum in
electronic excitation energy $\omega$, the same physical quantity as $E_e$
in Section~\ref{sec:dme}.

\subsection{\SrCdSb{} case: explicitly out of scope}

Unlike the DM-electron and dark-photon channels, the Migdal effect for
\SrCdSb{} is \textbf{not currently computed} in this framework — and should
not be reported as a projection. The reason is structural, not a missing
config:
\begin{itemize}[leftmargin=*]
  \item $I(\omega)$ requires the material's electron ELF
  $\Im[-1/\varepsilon(\omega,k)]$ over a \textbf{finite-$k$ grid}, exactly
  the same input DarkELF needs for absorption. DarkELF ships this for Si
  (\texttt{Si\_mermin.dat}) but has \textbf{no equivalent file for \SrCdSb{}}.
  \item Unlike the dark-photon channel, there is no simple Drude substitute
  here: $I(\omega)$ is a \textbf{$k$-integrated} quantity ($0\to k_\mathrm{max}$),
  so a frequency-only model (Drude $\varepsilon(\omega)$ at $k=0$) cannot
  stand in — the Migdal rate genuinely needs the momentum-dependent ELF,
  which is precisely the finite-$q$ information the Drude bracketing
  strategy was built to avoid needing.
  \item The scissors trick used for DM--electron scattering
  (Section~\ref{sec:dme-hypmat}) shifts \emph{Si's} band structure; it does
  not, by itself, produce $Z_\mathrm{ion}(k)$ or a finite-$k$ ELF file in the
  format DarkELF's Migdal module expects.
\end{itemize}

\paragraph{What would be required to add this channel:} a DarkELF-format
ELF grid $\Im[-1/\varepsilon(\omega,k)]$ for \SrCdSb{} over the relevant
$(\omega,k)$ range, plus an effective ionic charge $Z_\mathrm{ion}(k)$ for
the constituent atoms (Sr, Cd, Sb). Both require either full ab-initio
calculation or, at minimum, a material-specific momentum-dependent model —
beyond what the Drude/$\sigma_\mathrm{DC}$ bracketing approach can supply.
This is recorded as explicitly out of scope in
\texttt{Migdal\_Integration\_Plan.md} and remains so here.
